\chapter{svnversion Reference---Subversion Working Copy Version
Info}\label{svn.ref.svnversion}

\section{svnversion}\label{svn.ref.svnversion.re}

\emph{Summarize the local revision(s) of a working copy.}

\textbf{Synopsis}

\texttt{svnversion\ {[}OPTIONS{]}\ {[}WC\_PATH\ {[}TRAIL\_URL{]}{]}}

\subsection{Description}\label{svn.ref.svnversion.re.desc}

\texttt{svnversion} is a program for summarizing the revision mixture of
a working copy. The resultant revision number, or revision range, is
written to standard output.

It\textquotesingle s common to use this output in your build process
when defining the version number of your program.

\textless TRAIL\_URL\textgreater, if present, is the trailing portion of
the URL used to determine whether \textless WC\_PATH\textgreater{}
itself is switched (detection of switches within
\textless WC\_PATH\textgreater{} does not rely on
\textless TRAIL\_URL\textgreater).

When \textless WC\_PATH\textgreater{} is not defined, the current
directory will be used as the working copy path.
\textless TRAIL\_URL\textgreater{} cannot be defined if
\textless WC\_PATH\textgreater{} is not explicitly given.

\subsection{Options}\label{svn.ref.svnversion.re.sw}

Like \texttt{svnserve}, \texttt{svnversion} has no subcommands---only
options:

\begin{description}
\item[\protect\phantomsection\label{svn.ref.svnversion.sw.no_newline}{}\texttt{-\/-no-newline}
(\texttt{-n})]
Omits the usual trailing newline from the output.
\item[\protect\phantomsection\label{svn.ref.svnversion.sw.committed}{}\texttt{-\/-committed}
(\texttt{-c})]
Uses the last-changed revisions rather than the current (i.e., highest
locally available) revisions.
\item[\protect\phantomsection\label{svn.ref.svnversion.sw.help}{}\texttt{-\/-help}
(\texttt{-h})]
Prints a help summary.
\item[\protect\phantomsection\label{svn.ref.svnversion.sw.quiet}{}\texttt{-\/-quiet}
(\texttt{-q})]
Requests that the program print only essential information while
performing an operation.
\item[\protect\phantomsection\label{svn.ref.svnversion.sw.version}{}\texttt{-\/-version}]
Prints the version of \texttt{svnversion} and exit with no error.
\end{description}

\subsection{Examples}\label{svn.ref.svnversion.re.examples}

If the working copy is all at the same revision (e.g., immediately after
an update), then that revision is printed out:

\begin{verbatim}
$ svnversion
4168
\end{verbatim}

You can add \textless TRAIL\_URL\textgreater{} to make sure the working
copy is not switched from what you expect. Note that the
\textless WC\_PATH\textgreater{} is required in this command:

\begin{verbatim}
$ svnversion . /var/svn/trunk
4168
\end{verbatim}

For a mixed-revision working copy, the range of revisions present is
printed:

\begin{verbatim}
$ svnversion
4123:4168
\end{verbatim}

If the working copy contains modifications, a trailing
\textquotesingle{}\texttt{M}\textquotesingle{} is added:

\begin{verbatim}
$ svnversion
4168M
\end{verbatim}

If the working copy is switched, a trailing
\textquotesingle{}\texttt{S}\textquotesingle{} is added:

\begin{verbatim}
$ svnversion
4168S
\end{verbatim}

\texttt{svnversion} will also inform you if the target working copy is
sparsely populated (see \hyperref[svn.advanced.sparsedirs]{Sparse
Directories}) by attaching the
\textquotesingle{}\texttt{P}\textquotesingle{} code to its output:

\begin{verbatim}
$ svnversion
4168P
\end{verbatim}

Thus, here is a mixed-revision, sparsely populated and switched working
copy containing some local modifications:

\begin{verbatim}
$ svnversion
4123:4168MSP
\end{verbatim}

